\documentclass{article}

\usepackage{amsmath}

\title{Ch. 26 HW \#3}
\author{Principles of Macroeconomics}
\date{\today}

\begin{document}

\maketitle

\section*{Ch. 26 Problems and Applications}

\begin{enumerate}

\setcounter{enumi}{1}

\item

	\begin{enumerate}
	
	\item In the short run, it would increase economic growth because the faster growth of the capital stock would lead to faster growth in productivity. In the long run, the benefits from additional capital would fade out due to the diminishing marginal product of capital. As a result, the growth rate would return to its original levels in the long run.
	
	\item Consumers today might be hurt because they're asked to reduce their consumption. Consumers in the future would benefit because they would be richer and therefore capable of consuming more if they wished.
	
	\end{enumerate}

\item

	\begin{enumerate}
	
	\item Private consumption includes buying food, clothes, and train tickets. Private investment includes buying a new home, buying a laptop for a new employee, and building a factory.
	
	In my opinion, education spending is investment in human capital because you're increasing your skill level so that you can produce more in the future.
	
	\item Public consumption includes defense spending and spending on public sanitation. Public investment includes spending on transportation infrastructure and buildings for public employees. 
	
	I would consider health programs for anyone in or entering the labor force an investment, and I would consider health programs for anyone permanently out of the labor force as consumption. For people in the labor force, health programs increase their future capacity to work and therefore produce. For people out of the labor force, health programs increase their quality of life but not necessarily the productive capacity of the economy. 
	
	\end{enumerate}

\item The opportunity cost of investing in capital is consumption. Increasing investment necessarily requires a decrease in consumption, so society has to sacrifice consumption today in order to increase its productive capacity in the future.

Yes, it’s possible to overinvest in capital. At a certain point, the value of increasing future income isn’t worth giving up additional consumption today. 

The opportunity cost of investing in human capital is often forgone wages. Typically, investing in human capital involves enrolling in school or some kind of training program. Any time spent in school or a training program is time that could be spent working and earning a wage.

Yes, it’s possible to overinvest in human capital. At a certain point, the value of increasing our future earnings isn’t worth giving up additional earnings in the present.

\item

	\begin{enumerate}
	
	\item The investment allowed us to increase our capital stock and therefore our productive capacity. The fact that it was foreign allowed us to do so without reducing our current consumption levels. If we hadn’t received the investment, we would’ve consumed the same amount in the present, but we wouldn’t have been able to produce as much in the future.
	
	\item If Americans had made the investment themselves, then when future production increased, all of the extra future income generated would have accrued to Americans.
	
	\end{enumerate}

\end{enumerate}

\section*{Ch. 27 Quick Quiz}

\begin{enumerate}

\item (d) -- A bond is essentially a loan, and a stock is an ownership share, so Clarke is a bondholder, and Jimmy is a stockholder.

\item (d) -- The higher the credit risk, the higher the interest rate.

\item (b) -- Mutual funds tend to whole multiple assets, so they provide diversification.

\end{enumerate}

\section*{Ch. 27 Questions for Review}

\begin{enumerate}

\item The role of the financial system is to channel money from savers to investors.

One financial market is the bond market where people essentially buy IOUs from firms and governments. Another is the stock market where people buy ownership shares of firms. 

One financial intermediary is a bank which takes deposits from savers and uses them to make loans to investors. Another is a mutual fund which pools money from savers and uses the proceeds to buy a portfolio of stocks and bonds. 

\item Diversification allows investors to decrease the risk of their portfolio by holding a small share of a lot of companies. Mutual funds make diversification easier because investors can buy a share of a mutual fund and own a small part of the mutual fund's entire portfolio.

\end{enumerate}

\end{document}